\part{Entgegengesetzte Interessen}

\chapter{Was ein Spiel ist}

\section{Der Leitfall}

\begin{merksatz}
\boxtitel{Zwei Notarztstandorte}
Neustadt hat zwei Häuser: das Klinikum und das Marienhospital. Beide betreiben einen Notarztwagen und wählen dessen Standort --- Nord, Mitte oder Süd.

\begin{center}
\begin{minipage}{0.9\textwidth}
\textbf{Wer näher am Einsatzort steht, wird gerufen.} Der Anteil, den das eine Haus gewinnt, verliert das andere.
\end{minipage}
\end{center}

\schritt{Die Zahlen}
Anteil der Einsätze, die das Klinikum erreicht, in Prozent:

\begin{center}
{\small
\renewcommand{\arraystretch}{1.55}
\begin{tabular}{@{}P{2.2cm}P{1.6cm}P{1.6cm}P{1.6cm}@{}}
\toprule
\textbf{Klinikum} & \textbf{M.: N} & \textbf{M.: M} & \textbf{M.: S}\\
\midrule
Nord & $55$ & $40$ & $45$\\
Mitte & $60$ & $\mathbf{50}$ & $52$\\
Süd & $48$ & $38$ & $42$\\
\bottomrule
\end{tabular}}
\end{center}

\begin{center}
\begin{minipage}{0.88\textwidth}
\kursiv{Nur \kursiv{eine} Zahl je Feld genügt, weil das Marienhospital genau den Rest erhält. Bei $50$ für das Klinikum bleiben dem Marienhospital $50$. \textbf{Solche Spiele heißen Nullsummenspiele} --- nach einer Verschiebung der Skala summieren sich die Auszahlungen zu null.}
\end{minipage}
\end{center}
\end{merksatz}

\section{Die Bestandteile}

\begin{definition}
\boxtitelm{Definition 1.1 --- Spiel in Normalform}
Ein \Begriff{Spiel in Normalform} besteht aus
\begin{center}
{\small
\renewcommand{\arraystretch}{1.5}
\begin{tabular}{@{}P{2.4cm}P{6.2cm}@{}}
\toprule
\textbf{Spielern} & einer endlichen Menge $\{1,\ldots,n\}$\\
\textbf{Strategien} & für jeden Spieler $i$ einer Menge $\Str_i$\\
\textbf{Auszahlungen} & für jeden Spieler einer Funktion $\Aus_i: \Str_1\times\cdots\times\Str_n \to \R$\\
\bottomrule
\end{tabular}}
\end{center}
\end{definition}

\begin{merksatz}
\boxtitel{Was in dieser Definition steckt --- und was nicht}
\begin{center}
{\small
\renewcommand{\arraystretch}{1.52}
\begin{tabular}{@{}P{2.4cm}P{6.2cm}@{}}
\toprule
\textbf{Gleichzeitig} & Alle wählen, ohne die Wahl der anderen zu kennen. Kapitel~7 lockert das.\\
\textbf{Vollständig bekannt} & Jeder kennt alle Strategiemengen und alle Auszahlungen --- auch die der anderen.\\
\textbf{Eine Zahl je Ausgang} & Alles, was einem Spieler wichtig ist, steckt in einer einzigen Zahl.\\
\bottomrule
\end{tabular}}
\end{center}

\begin{center}
\begin{minipage}{0.9\textwidth}
\textbf{Die zweite Zeile ist eine starke Annahme.} Sie besagt nicht nur, dass jeder die Auszahlungen kennt, sondern auch, dass jeder weiß, dass jeder sie kennt --- und so fort.
\end{minipage}
\end{center}

\begin{center}
\begin{minipage}{0.88\textwidth}
\kursiv{Man nennt das \Begriff{allgemeines Wissen}. Ohne diese unendliche Kette bricht ein Teil der Schlüsse zusammen, die in Abschnitt 1.4 gezogen werden --- denn dort schließt Spieler A daraus, was B aus dem schließt, was A tut. \textbf{Die Annahme ist selten erfüllt und fast nie überprüft.}}
\end{minipage}
\end{center}

\schritt{Zur dritten Zeile}
Die Auszahlung ist kein Geldbetrag, sondern ein \Begriff{Nutzen}: eine Zahl, die die Vorlieben des Spielers vollständig ordnet.

\begin{center}
\begin{minipage}{0.88\textwidth}
\kursiv{Für das Klinikum steckt darin der Erlös, aber auch der Ruf, die Auslastung, die Zufriedenheit der Mitarbeiter --- alles in einer Zahl verrechnet. \textbf{Dass sich das überhaupt tun lässt, ist ein eigener Satz}, den von Neumann und Morgenstern 1944 bewiesen und den Kapitel~4 behandelt. Er verlangt Annahmen, die man sich ansehen sollte, bevor man ihm folgt.}
\end{minipage}
\end{center}
\end{merksatz}

\begin{definition}
\boxtitelm{Definition 1.2 --- Nullsummenspiel}
Ein Zweipersonenspiel heißt \Begriff{Nullsummenspiel}, wenn
\[ \Aus_1(s_1,s_2) + \Aus_2(s_1,s_2) = 0 \qquad\text{für alle } (s_1,s_2) . \]
\end{definition}

\begin{merksatz}
\boxtitel{Warum eine Matrix genügt}
Ist $\Aus_2 = -\Aus_1$, so ist das Spiel durch die Matrix von $\Aus_1$ vollständig beschrieben.

\begin{center}
\begin{minipage}{0.9\textwidth}
\textbf{Der Zeilenspieler will den Eintrag groß machen, der Spaltenspieler klein.}
\end{minipage}
\end{center}

\begin{center}
\begin{minipage}{0.88\textwidth}
\kursiv{Im Leitfall summieren sich die Anteile zu $100$, nicht zu null --- aber eine Verschiebung um $50$ macht daraus ein Nullsummenspiel, ohne irgendeine Entscheidung zu ändern. \textbf{Was zählt, ist die Gegenläufigkeit, nicht die Zahl null.} Man spricht deshalb oft allgemeiner von \Begriff{strikt kompetitiven} Spielen.}
\end{minipage}
\end{center}
\end{merksatz}

\section{Sich absichern}

\begin{merksatz}
\boxtitel{Die vorsichtige Rechnung}
Das Klinikum fragt sich: Was ist mir garantiert, wenn ich Zeile $r$ wähle und das Marienhospital anschließend das Schlimmste tut?

\begin{center}
{\small
\renewcommand{\arraystretch}{1.55}
\begin{tabular}{@{}P{2.0cm}P{1.4cm}P{1.4cm}P{1.4cm}P{2.0cm}@{}}
\toprule
& \textbf{N} & \textbf{M} & \textbf{S} & \textbf{Garantie}\\
\midrule
Nord & $55$ & $40$ & $45$ & $40$\\
Mitte & $60$ & $50$ & $52$ & $\mathbf{50}$\\
Süd & $48$ & $38$ & $42$ & $38$\\
\bottomrule
\end{tabular}}
\end{center}

\[ \text{Maximin} = \max_r \min_c a_{rc} = 50 \]

\schritt{Die andere Seite}
Das Marienhospital fragt: Was riskiere ich höchstens, wenn ich Spalte $c$ wähle?

\begin{center}
{\small
\renewcommand{\arraystretch}{1.55}
\begin{tabular}{@{}P{2.8cm}P{1.3cm}P{1.3cm}P{1.3cm}@{}}
\toprule
& \textbf{N} & \textbf{M} & \textbf{S}\\
\midrule
Spaltenmaximum & $60$ & $\mathbf{50}$ & $52$\\
\bottomrule
\end{tabular}}
\end{center}

\[ \text{Minimax} = \min_c \max_r a_{rc} = 50 \]
\end{merksatz}

\begin{merksatz}
\boxtitel{Satz 1.3 --- Maximin liegt nie über Minimax}
Für jede Matrix gilt
\[ \max_r \min_c a_{rc} \ \le\ \min_c \max_r a_{rc} . \]

\schritt{Beweis}
Seien $r^{*}$ und $c^{*}$ beliebig. Dann gilt
\[ \min_c a_{r^{*}c} \ \le\ a_{r^{*}c^{*}} \ \le\ \max_r a_{rc^{*}} . \]

\kursiv{Erste Ungleichung:} Das Minimum über alle $c$ ist höchstens der Wert bei $c^{*}$.\\
\kursiv{Zweite:} Der Wert bei $r^{*}$ ist höchstens das Maximum über alle $r$.

Die linke Seite hängt nicht von $c^{*}$ ab, die rechte nicht von $r^{*}$. Also gilt die Ungleichung auch für das Maximum links und das Minimum rechts. $\blacksquare$

\begin{center}
\begin{minipage}{0.88\textwidth}
\kursiv{Der Beweis ist drei Zeilen lang und benutzt nichts als die Definition von Maximum und Minimum. \textbf{Er gilt für jede Matrix, ohne jede Voraussetzung} --- und er sagt, dass die vorsichtige Garantie des einen nie besser ist als das schlimmste Zugeständnis des anderen.}
\end{minipage}
\end{center}
\end{merksatz}

\begin{definition}
\boxtitelm{Definition 1.4 --- Sattelpunkt}
Ein Feld $(r^{*},c^{*})$ heißt \Begriff{Sattelpunkt}, wenn es zugleich das Minimum seiner Zeile und das Maximum seiner Spalte ist.
\end{definition}

\begin{merksatz}
\boxtitel{Satz 1.5 --- Sattelpunkt heißt Gleichheit}
Ein Sattelpunkt existiert genau dann, wenn Maximin und Minimax übereinstimmen. Sein Wert ist dann dieser gemeinsame Wert.

\schritt{Beweis, Richtung eins}
Sei $(r^{*},c^{*})$ ein Sattelpunkt mit Wert $v = a_{r^{*}c^{*}}$.
\begin{align*}
\text{Maximin} &\ \ge\ \min_c a_{r^{*}c} = v && \text{($v$ ist Zeilenminimum)}\\
\text{Minimax} &\ \le\ \max_r a_{rc^{*}} = v && \text{($v$ ist Spaltenmaximum)}
\end{align*}
Zusammen mit Satz 1.3:
\[ v \le \text{Maximin} \le \text{Minimax} \le v \]
Also sind alle drei gleich. \checkmark

\schritt{Richtung zwei}
Gilt Maximin $=$ Minimax $= v$, so gibt es ein $r^{*}$ mit $\min_c a_{r^{*}c} = v$ und ein $c^{*}$ mit $\max_r a_{rc^{*}} = v$.

Dann ist
\[ v = \min_c a_{r^{*}c} \ \le\ a_{r^{*}c^{*}} \ \le\ \max_r a_{rc^{*}} = v , \]
also $a_{r^{*}c^{*}} = v$ --- und das Feld ist Zeilenminimum und Spaltenmaximum zugleich. $\blacksquare$
\end{merksatz}

\begin{merksatz}
\boxtitel{Der Leitfall hat einen}
\[ \text{Maximin} = \text{Minimax} = 50 \]

\begin{center}
{\small
\renewcommand{\arraystretch}{1.55}
\begin{tabular}{@{}P{2.2cm}P{1.2cm}P{1.2cm}P{1.2cm}P{1.9cm}@{}}
\toprule
& \textbf{N} & \textbf{M} & \textbf{S} & \textbf{Zeilenmin}\\
\midrule
Nord & $55$ & $40$ & $45$ & $40$\\
Mitte & $60$ & \cellcolor{akzent!18}$\mathbf{50}$ & $52$ & $\mathbf{50}$\\
Süd & $48$ & $38$ & $42$ & $38$\\
\midrule
\textbf{Spaltenmax} & $60$ & $\mathbf{50}$ & $52$ & \\
\bottomrule
\end{tabular}}
\end{center}

\begin{center}
\begin{minipage}{0.9\textwidth}
\textbf{Beide Häuser wählen Mitte, das Klinikum erreicht $50$ Prozent, und keines hat einen Grund, davon abzuweichen.}
\end{minipage}
\end{center}

\schritt{Was das bedeutet}
\begin{center}
{\small
\renewcommand{\arraystretch}{1.5}
\begin{tabular}{@{}P{2.6cm}P{6.0cm}@{}}
\toprule
\textbf{Kein Geheimnis nötig} & Das Klinikum könnte seine Wahl vorher bekanntgeben --- am Ergebnis änderte sich nichts.\\
\textbf{Kein Bedauern} & Nachdem beide gewählt haben, bereut keiner.\\
\textbf{Ein Wert} & Die Zahl $50$ beschreibt das Spiel vollständig.\\
\bottomrule
\end{tabular}}
\end{center}

\begin{center}
\begin{minipage}{0.88\textwidth}
\kursiv{Die erste Zeile ist das eigentliche Kennzeichen. In einem Spiel mit Sattelpunkt ist Verschwiegenheit wertlos --- und in einem ohne, wie wir in Kapitel~2 sehen werden, ist sie alles. \textbf{Der Sattelpunkt ist genau die Lage, in der Voraussagbarkeit nicht schadet.}}
\end{minipage}
\end{center}
\end{merksatz}

\section{Dominanz}

\begin{definition}
\boxtitelm{Definition 1.6}
Eine Strategie $s$ \Begriff{dominiert} eine Strategie $s'$ \Begriff{strikt}, wenn sie gegen \kursiv{jede} Wahl der Gegenseite echt besser abschneidet.
\end{definition}

\begin{merksatz}
\boxtitel{Der Leitfall löst sich auch so}
\schritt{Schritt 1 --- Für das Klinikum}
Zeile \enquote{Mitte} gegen Zeile \enquote{Nord}:
\[ 60 > 55, \qquad 50 > 40, \qquad 52 > 45 \]
\textbf{Mitte dominiert Nord strikt.} Ebenso Süd:
\[ 60 > 48, \qquad 50 > 38, \qquad 52 > 42 \]

\schritt{Schritt 2 --- Für das Marienhospital}
Es will kleine Werte. Spalte \enquote{M} gegen die anderen:
\[ 40 < 55, \quad 50 < 60, \quad 38 < 48 \qquad\text{(gegen N)} \]
\[ 40 < 45, \quad 50 < 52, \quad 38 < 42 \qquad\text{(gegen S)} \]
\textbf{Mitte dominiert auch hier strikt.}

\schritt{Schritt 3 --- Das Ergebnis}
\begin{center}
\begin{minipage}{0.9\textwidth}
\textbf{Beide Seiten haben eine Strategie, die immer besser ist.} Man braucht den Sattelpunkt gar nicht zu suchen.
\end{minipage}
\end{center}

\begin{center}
\begin{minipage}{0.88\textwidth}
\kursiv{Das ist hier ein glücklicher Sonderfall. \textbf{Dominanz ist die stärkste Art von Argument in diesem Band}, weil sie ohne jede Annahme über das Verhalten des Gegenübers auskommt: Man muss nicht unterstellen, dass er vernünftig handelt, sondern nur, dass man selbst es tut.}
\end{minipage}
\end{center}
\end{merksatz}

\begin{merksatz}
\boxtitel{Schrittweise streichen}
Meist dominiert nicht eine Strategie alle anderen, sondern nur einzelne. Dann streicht man und schaut erneut.

\schritt{Ein Beispiel}
\begin{center}
{\small
\renewcommand{\arraystretch}{1.5}
\begin{tabular}{@{}P{1.6cm}P{1.4cm}P{1.4cm}P{1.4cm}@{}}
\toprule
& \textbf{L} & \textbf{M} & \textbf{R}\\
\midrule
\textbf{O} & $4$ & $3$ & $7$\\
\textbf{U} & $2$ & $6$ & $5$\\
\bottomrule
\end{tabular}}
\end{center}

\kursiv{Für den Spaltenspieler:} L $(4;2)$ gegen R $(7;5)$ --- L ist in beiden Zeilen kleiner. \textbf{L dominiert R.} Streichen:

\begin{center}
{\small
\renewcommand{\arraystretch}{1.5}
\begin{tabular}{@{}P{1.6cm}P{1.4cm}P{1.4cm}@{}}
\toprule
& \textbf{L} & \textbf{M}\\
\midrule
\textbf{O} & $4$ & $3$\\
\textbf{U} & $2$ & $6$\\
\bottomrule
\end{tabular}}
\end{center}

\kursiv{Jetzt:} Keine weitere Dominanz. Zeilenminima $3$ und $2$, Spaltenmaxima $4$ und $6$.
\[ \text{Maximin} = 3 \ <\ 4 = \text{Minimax} \]

\begin{center}
\begin{minipage}{0.9\textwidth}
\textbf{Kein Sattelpunkt.} Die Dominanz half, reichte aber nicht.
\end{minipage}
\end{center}

\begin{center}
\begin{minipage}{0.88\textwidth}
\kursiv{Das ist der Normalfall, und er ist der Grund für Kapitel~2 und 3. Die Lücke zwischen $3$ und $4$ ist nicht zu schließen, solange man sich auf feste Wahlen beschränkt. \textbf{Man muss würfeln} --- und der Satz, der das rechtfertigt, ist der wichtigste dieses Teils.}
\end{minipage}
\end{center}
\end{merksatz}

\begin{warnung}
\boxtitelr{Strikt oder schwach --- ein Unterschied mit Folgen}
Dominiert $s$ die Strategie $s'$ nur \kursiv{schwach} --- also nirgends schlechter, irgendwo besser ---, so kann das Streichen Lösungen vernichten.

\begin{center}
{\small
\renewcommand{\arraystretch}{1.5}
\begin{tabular}{@{}P{1.6cm}P{1.4cm}P{1.4cm}@{}}
\toprule
& \textbf{L} & \textbf{R}\\
\midrule
\textbf{O} & $1$ & $1$\\
\textbf{U} & $0$ & $2$\\
\bottomrule
\end{tabular}}
\end{center}

\kursiv{Für den Spaltenspieler} dominiert L die Strategie R schwach: $1 \le 1$ und $0 < 2$.

\begin{center}
\begin{minipage}{0.9\textwidth}
\textbf{Streicht man R, so bleibt die Lösung $(\text{O},\text{L})$ mit Wert $1$.} Das Feld $(\text{O},\text{R})$ mit demselben Wert ist ebenfalls ein Sattelpunkt --- und wurde weggeworfen.
\end{minipage}
\end{center}

\begin{center}
\begin{minipage}{0.88\textwidth}
\kursiv{Beim Nullsummenspiel ist der Schaden gering, weil der Wert derselbe bleibt. \textbf{Bei den Spielen aus Teil~II ist er erheblich}: Dort kann schwache Dominanzelimination Gleichgewichte mit ganz verschiedenen Auszahlungen beseitigen, und sogar die \kursiv{Reihenfolge} des Streichens beeinflusst das Ergebnis. Strikte Dominanz hat diese Schwäche nicht.}
\end{minipage}
\end{center}
\end{warnung}

\section{Der Bruchfall}

\begin{warnung}
\boxtitelr{Was die Matrix verschweigt}
\begin{center}
{\small
\renewcommand{\arraystretch}{1.52}
\begin{tabular}{@{}P{2.4cm}P{6.2cm}@{}}
\toprule
\textbf{Woher die Zahlen} & Die Auszahlungen sind Schätzungen. Im Leitfall beruhen sie auf Fahrzeiten, Bevölkerungsdichte und Annahmen über die Disposition.\\
\textbf{Wer die Spieler sind} & Ein Haus ist kein Spieler --- Geschäftsführung, Chefärzte und Träger haben verschiedene Ziele.\\
\textbf{Was sonst geschieht} & Die Leitstelle disponiert nach eigenen Regeln. Sie ist ein dritter Beteiligter, der nicht in der Matrix steht.\\
\bottomrule
\end{tabular}}
\end{center}

\schritt{Zur zweiten Zeile}
\begin{center}
\begin{minipage}{0.88\textwidth}
\kursiv{Die Annahme, eine Organisation handle wie eine Person mit einer einzigen Zielfunktion, ist bequem und meist falsch. \textbf{Sie ist zugleich unvermeidlich}, wenn man überhaupt rechnen will --- und deshalb gehört zu jeder spieltheoretischen Analyse die Frage, wie empfindlich das Ergebnis auf diese Vereinfachung reagiert.}
\end{minipage}
\end{center}

\schritt{Und die Empfindlichkeit prüfen}
Im Leitfall liegt der Sattelpunkt bei $50$, die nächstbeste Zeile garantiert $40$.

\begin{center}
\begin{minipage}{0.9\textwidth}
\textbf{Der Abstand beträgt zehn Punkte.} Eine Schätzung, die um zwei oder drei danebenliegt, ändert nichts --- eine, die um zwölf danebenliegt, alles.
\end{minipage}
\end{center}

\begin{center}
\begin{minipage}{0.88\textwidth}
\kursiv{Diese Rechnung ist die Entsprechung der Schattenpreise und Gültigkeitsbereiche aus Band 19, Kapitel~5. \textbf{Sie sollte zu jeder Lösung dazugehören} und fehlt in den meisten Anwendungen --- was die erste Warnung des Vorworts erklärt.}
\end{minipage}
\end{center}
\end{warnung}

\section{Übungsaufgaben}

\begin{aufgabe}
\boxtitel{Aufgabe 1.1 --- Sattelpunkte finden}
Bestimmen Sie für jede Matrix Maximin, Minimax und gegebenenfalls den Sattelpunkt.
\begin{center}
{\small
\renewcommand{\arraystretch}{1.44}
\begin{tabular}{@{}P{1.0cm}P{2.6cm}P{1.0cm}P{2.6cm}@{}}
\toprule
(a) & $\begin{pmatrix}3&5\\2&4\end{pmatrix}$ & (b) & $\begin{pmatrix}1&4\\3&2\end{pmatrix}$\\[12pt]
(c) & $\begin{pmatrix}6&2&8\\4&5&7\\3&1&9\end{pmatrix}$ & (d) & $\begin{pmatrix}2&2\\2&2\end{pmatrix}$\\
\bottomrule
\end{tabular}}
\end{center}

\loesung
\schritt{(a)}
\begin{center}
\renewcommand{\arraystretch}{1.46}
\begin{tabular}{@{}P{2.2cm}P{1.1cm}P{1.1cm}P{1.9cm}@{}}
\toprule
& $c_1$ & $c_2$ & \textbf{Zeilenmin}\\
\midrule
$r_1$ & $3$ & $5$ & $\mathbf{3}$\\
$r_2$ & $2$ & $4$ & $2$\\
\midrule
\textbf{Spaltenmax} & $\mathbf{3}$ & $5$ & \\
\bottomrule
\end{tabular}
\end{center}

\[ \text{Maximin} = 3 = \text{Minimax} \quad\Longrightarrow\quad \text{Sattelpunkt bei } (r_1,c_1), \ v = 3 \]

\kursiv{Probe:} $3$ ist Minimum der Zeile $(3;5)$ \checkmark und Maximum der Spalte $(3;2)$ \checkmark

\schritt{(b)}
\begin{center}
\renewcommand{\arraystretch}{1.46}
\begin{tabular}{@{}P{2.2cm}P{1.1cm}P{1.1cm}P{1.9cm}@{}}
\toprule
& $c_1$ & $c_2$ & \textbf{Zeilenmin}\\
\midrule
$r_1$ & $1$ & $4$ & $1$\\
$r_2$ & $3$ & $2$ & $\mathbf{2}$\\
\midrule
\textbf{Spaltenmax} & $\mathbf{3}$ & $4$ & \\
\bottomrule
\end{tabular}
\end{center}

\[ \text{Maximin} = 2 \ <\ 3 = \text{Minimax} \]

\textcolor{akzent}{\textbf{Kein Sattelpunkt.}} Die Lücke beträgt $1$.

\schritt{(c)}
\begin{center}
\renewcommand{\arraystretch}{1.46}
\begin{tabular}{@{}P{2.2cm}P{1.0cm}P{1.0cm}P{1.0cm}P{1.9cm}@{}}
\toprule
& $c_1$ & $c_2$ & $c_3$ & \textbf{Zeilenmin}\\
\midrule
$r_1$ & $6$ & $2$ & $8$ & $2$\\
$r_2$ & $4$ & $5$ & $7$ & $\mathbf{4}$\\
$r_3$ & $3$ & $1$ & $9$ & $1$\\
\midrule
\textbf{Spaltenmax} & $6$ & $\mathbf{5}$ & $9$ & \\
\bottomrule
\end{tabular}
\end{center}

\[ \text{Maximin} = 4 \ <\ 5 = \text{Minimax} \]

\textcolor{akzent}{\textbf{Kein Sattelpunkt.}}

\schritt{(d)}
Alle Einträge gleich: Maximin $=$ Minimax $= 2$.

\begin{center}
\begin{minipage}{0.9\textwidth}
\textbf{Alle vier Felder sind Sattelpunkte.} Sattelpunkte sind nicht eindeutig.
\end{minipage}
\end{center}

\schritt{Die Übersicht}
\begin{center}
{\small
\renewcommand{\arraystretch}{1.48}
\begin{tabular}{@{}P{1.0cm}P{1.8cm}P{1.8cm}P{2.2cm}@{}}
\toprule
& \textbf{Maximin} & \textbf{Minimax} & \textbf{Sattelpunkt}\\
\midrule
(a) & $3$ & $3$ & \textcolor{gruen}{ja, eindeutig}\\
(b) & $2$ & $3$ & \textcolor{akzent}{nein}\\
(c) & $4$ & $5$ & \textcolor{akzent}{nein}\\
(d) & $2$ & $2$ & \textcolor{gruen}{ja, vier}\\
\bottomrule
\end{tabular}}
\end{center}

\begin{center}
\begin{minipage}{0.88\textwidth}
\kursiv{Zwei von vier Matrizen haben keinen Sattelpunkt --- und das ist typisch. \textbf{Bei zufällig gewählten $2\times2$-Matrizen liegt der Anteil bei zwei Dritteln}, bei größeren Matrizen sinkt er rasch. Kapitel~2 zeigt, dass sich die Lücke stets schließen lässt, wenn man Zufall zulässt.}
\end{minipage}
\end{center}
\end{aufgabe}

\begin{aufgabe}
\boxtitel{Aufgabe 1.2 --- Dominanz anwenden}
Vereinfachen Sie durch schrittweises Streichen strikt dominierter Strategien:
\[ \begin{pmatrix} 5 & 3 & 6 & 4\\ 7 & 4 & 8 & 5\\ 2 & 1 & 3 & 2 \end{pmatrix} \]

\loesung
\schritt{Schritt 1 --- Zeilen prüfen (Zeilenspieler maximiert)}
\kursiv{Zeile 2 gegen Zeile 1:} $7>5$, $4>3$, $8>6$, $5>4$ --- \textbf{Zeile 2 dominiert Zeile 1 strikt.}

\kursiv{Zeile 2 gegen Zeile 3:} $7>2$, $4>1$, $8>3$, $5>2$ --- \textbf{ebenfalls.}

Es bleibt nur Zeile 2:
\[ \begin{pmatrix} 7 & 4 & 8 & 5 \end{pmatrix} \]

\schritt{Schritt 2 --- Spalten prüfen (Spaltenspieler minimiert)}
\kursiv{Spalte 2} hat den kleinsten Wert $4$; sie dominiert alle anderen strikt:
\[ 4 < 7, \qquad 4 < 8, \qquad 4 < 5 \]

\schritt{Schritt 3 --- Das Ergebnis}
\[ \text{Lösung: } (r_2,c_2), \qquad v = 4 \]

\schritt{Schritt 4 --- Die Probe über Sattelpunkte}
\begin{center}
\renewcommand{\arraystretch}{1.46}
\begin{tabular}{@{}P{2.2cm}P{0.85cm}P{0.85cm}P{0.85cm}P{0.85cm}P{1.9cm}@{}}
\toprule
& $c_1$ & $c_2$ & $c_3$ & $c_4$ & \textbf{Zeilenmin}\\
\midrule
$r_1$ & $5$ & $3$ & $6$ & $4$ & $3$\\
$r_2$ & $7$ & $4$ & $8$ & $5$ & $\mathbf{4}$\\
$r_3$ & $2$ & $1$ & $3$ & $2$ & $1$\\
\midrule
\textbf{Spaltenmax} & $7$ & $\mathbf{4}$ & $8$ & $5$ & \\
\bottomrule
\end{tabular}
\end{center}

\[ \text{Maximin} = \text{Minimax} = 4 \quad \checkmark \]

\schritt{Schritt 5 --- Eine Beobachtung zur Reihenfolge}
Hätte man mit den Spalten begonnen, so wäre Spalte 2 ebenfalls sofort übrig geblieben --- und dann Zeile 2.

\begin{center}
\begin{minipage}{0.88\textwidth}
\kursiv{Das ist kein Zufall: \textbf{Bei strikter Dominanz ist das Ergebnis von der Reihenfolge unabhängig}, und die Menge der überlebenden Strategien ist eindeutig bestimmt. Bei schwacher Dominanz gilt das nicht, wie der Warnkasten in Abschnitt 1.4 zeigte. Der Beweis der Unabhängigkeit steht in Osborne und Rubinstein\zit{19}.}
\end{minipage}
\end{center}
\end{aufgabe}

\begin{aufgabe}
\boxtitel{Aufgabe 1.3 --- Die Empfindlichkeit prüfen}
Im Leitfall wurde der Wert $50$ für das Feld (Mitte, Mitte) geschätzt. Um wie viel darf die Schätzung danebenliegen, ohne dass sich die Lösung ändert?

\loesung
\schritt{Schritt 1 --- Was der Sattelpunkt verlangt}
Der Wert $a$ im Feld (Mitte, Mitte) muss
\begin{center}
{\small
\renewcommand{\arraystretch}{1.5}
\begin{tabular}{@{}P{3.4cm}P{5.0cm}@{}}
\toprule
Zeilenminimum sein & $a \le 60$ und $a \le 52$\\
Spalten\-maximum sein & $a \ge 40$ und $a \ge 38$\\
\bottomrule
\end{tabular}}
\end{center}

\schritt{Schritt 2 --- Der Bereich}
\[ 40 \ \le\ a \ \le\ 52 \]

\begin{center}
\begin{minipage}{0.9\textwidth}
\textbf{Die Schätzung darf um $10$ nach unten und um $2$ nach oben danebenliegen.}
\end{minipage}
\end{center}

\schritt{Schritt 3 --- Was außerhalb geschieht}
\kursiv{Bei $a = 55$:} Das Feld ist kein Zeilenminimum mehr ($55 > 52$).
\begin{center}
\renewcommand{\arraystretch}{1.46}
\begin{tabular}{@{}P{1.8cm}P{1.2cm}P{1.2cm}P{1.2cm}P{1.6cm}@{}}
\toprule
& N & M & S & \textbf{Min}\\
\midrule
Nord & $55$ & $40$ & $45$ & $40$\\
Mitte & $60$ & $55$ & $52$ & $\mathbf{52}$\\
Süd & $48$ & $38$ & $42$ & $38$\\
\midrule
\textbf{Max} & $60$ & $55$ & $\mathbf{52}$ & \\
\bottomrule
\end{tabular}
\end{center}

Maximin $=$ Minimax $= 52$ --- der Sattelpunkt wandert nach (Mitte, Süd).

\kursiv{Bei $a = 35$:} Das Feld ist kein Spaltenmaximum mehr ($35 < 40$).
\begin{center}
\renewcommand{\arraystretch}{1.46}
\begin{tabular}{@{}P{1.8cm}P{1.2cm}P{1.2cm}P{1.2cm}P{1.6cm}@{}}
\toprule
& N & M & S & \textbf{Min}\\
\midrule
Nord & $55$ & $40$ & $45$ & $40$\\
Mitte & $60$ & $35$ & $52$ & $35$\\
Süd & $48$ & $38$ & $42$ & $38$\\
\midrule
\textbf{Max} & $60$ & $40$ & $52$ & \\
\bottomrule
\end{tabular}
\end{center}

Maximin $= 40$, Minimax $= 40$ --- Sattelpunkt bei (Nord, Mitte).

\schritt{Schritt 4 --- Die Übersicht}
\begin{center}
{\small
\renewcommand{\arraystretch}{1.48}
\begin{tabular}{@{}P{2.0cm}P{2.6cm}P{2.4cm}@{}}
\toprule
$a$ & \textbf{Sattelpunkt} & \textbf{Wert}\\
\midrule
$<40$ & (Nord, Mitte) & $40$\\
$40\ldots52$ & (Mitte, Mitte) & $a$\\
$>52$ & (Mitte, Süd) & $52$\\
\bottomrule
\end{tabular}}
\end{center}

\begin{center}
\begin{minipage}{0.9\textwidth}
\textbf{Der Wert des Spiels ist eine stetige, stückweise lineare Funktion der geschätzten Größe --- und er bleibt beschränkt.}
\end{minipage}
\end{center}

\schritt{Schritt 5 --- Die Deutung}
\begin{center}
\begin{minipage}{0.88\textwidth}
\kursiv{Selbst eine grob falsche Schätzung verändert den Wert nur zwischen $40$ und $52$. \textbf{Die Entscheidung des Klinikums bleibt in allen drei Fällen dieselbe: Mitte oder Nord}, und beide liegen nahe beieinander. Das ist ein beruhigender Befund --- und keineswegs selbstverständlich, wie Kapitel~5 an Beispielen zeigen wird, in denen eine Änderung um eine Einheit die Lösung vollständig umwirft.}
\end{minipage}
\end{center}
\end{aufgabe}

\begin{geschichte}
\boxtitelg{Ein Schachsatz und ein Beweis für sich}
\Person{Ernst Zermelo} (1871--1953) hielt 1912 auf einem Mathematikerkongress einen Vortrag über das Schachspiel.

\begin{center}
\begin{minipage}{0.9\textwidth}
\kursiv{Sein Ergebnis: Im Schach gibt es eine bestimmte Lage --- entweder kann Weiß einen Sieg erzwingen, oder Schwarz, oder beide können ein Remis erzwingen. \textbf{Welche der drei zutrifft, ist bis heute unbekannt} und wird es vermutlich bleiben. Der Satz sagt, dass eine Antwort existiert, nicht welche.}
\end{minipage}
\end{center}

\schritt{Die Methode}
Zermelos Beweis arbeitet rückwärts vom Spielende --- das Verfahren, das Kapitel~7 als Rückwärtsinduktion behandelt.

\begin{center}
\begin{minipage}{0.9\textwidth}
\kursiv{Es ist bemerkenswert, dass die Spieltheorie mit einem Satz über ein Spiel mit vollständiger Information beginnt, das sich gerade \kursiv{nicht} mit Matrizen behandeln lässt. \textbf{Die Normalform kam später und von einer anderen Seite.}}
\end{minipage}
\end{center}

\schritt{Borel}
\Person{Émile Borel} (1871--1956) untersuchte ab 1921 Spiele mit gemischten Strategien und formulierte die Frage nach dem Minimax-Satz.

\begin{center}
\begin{minipage}{0.9\textwidth}
\kursiv{Er hielt die allgemeine Antwort für zweifelhaft und vermutete, dass es Spiele ohne Wert gibt. \textbf{Borel lag falsch}, wie von Neumann 1928 zeigte --- was Borel später ohne Umschweife einräumte. Seine Arbeiten enthalten aber die Begriffe, die dafür nötig waren.}
\end{minipage}
\end{center}

\schritt{Und die Vorgeschichte}
\Person{James Waldegrave} beschrieb 1713 in einem Brief die gemischte Gleichgewichtsstrategie für ein Kartenspiel namens \kursiv{Le Her}.

\begin{center}
\begin{minipage}{0.9\textwidth}
\kursiv{Es ist die erste bekannte Minimax-Lösung überhaupt --- zweihundert Jahre vor Borel, in einem Briefwechsel über ein Gesellschaftsspiel. Der Brief wurde erst 1968 als solcher erkannt. \textbf{Es ist wieder das Muster aus Band 20: Eine Entdeckung, die niemand veröffentlicht, wirkt nicht.}}
\end{minipage}
\end{center}
\end{geschichte}

\begin{forschung}
\boxtitelp{Wo die Normalform an Grenzen stößt}
\begin{center}
{\small
\renewcommand{\arraystretch}{1.45}
\begin{tabular}{@{}P{2.8cm}P{5.4cm}@{}}
\toprule
\textbf{Größe} & Schon Schach hat mehr Strategien, als das Weltall Atome hat. Die Normalform ist dort kein brauchbares Werkzeug.\\
\textbf{Knapper Verstand} & Spieler rechnen nicht beliebig tief. Modelle mit begrenzter Denktiefe sagen Verhalten besser vorher.\\
\textbf{Unvollständiges Wissen} & Wer die Auszahlungen des anderen nicht kennt, spielt ein anderes Spiel.\\
\bottomrule
\end{tabular}}
\end{center}

\schritt{Zur zweiten Zeile}
In Versuchen wählen Menschen selten das Gleichgewicht, sondern denken ein bis zwei Schritte weit.

\begin{center}
\begin{minipage}{0.88\textwidth}
\kursiv{Modelle mit Denktiefen von null bis drei sagen beobachtetes Verhalten deutlich besser vorher als die Gleichgewichtsanalyse. \textbf{Die Dominanzargumente aus Abschnitt 1.4 sind davon am wenigsten betroffen} --- eine strikt dominierte Strategie wird auch von wenig rechnenden Menschen gemieden. Je mehr Schritte ein Argument braucht, desto schlechter beschreibt es die Wirklichkeit.}
\end{minipage}
\end{center}

\schritt{Zur dritten Zeile}
\Person{John Harsanyi} zeigte 1967, wie sich Spiele mit unvollständigem Wissen in Spiele mit vollständigem, aber unvollkommenem Wissen umformen lassen.

\begin{center}
\begin{minipage}{0.88\textwidth}
\kursiv{Der Kunstgriff: Man führt einen Zufallszug der Natur ein, der die Eigenschaften der Spieler bestimmt. Damit wird das Problem behandelbar --- um den Preis einer Annahme über die Wahrscheinlichkeiten, die niemand beobachtet. \textbf{Harsanyi erhielt dafür 1994 den Nobelpreis}, gemeinsam mit Nash und Selten.}
\end{minipage}
\end{center}
\end{forschung}

\begin{merksatz}
\boxtitel{Was Sie jetzt können}
\begin{center}
\renewcommand{\arraystretch}{1.5}
\begin{tabular}{@{}P{3.0cm}P{5.6cm}@{}}
\toprule
\textbf{Definition 1.1} & Spieler, Strategien, Auszahlungen\\
\textbf{Definition 1.2} & Nullsumme: eine Matrix genügt\\
\textbf{Satz 1.3} & Maximin $\le$ Minimax, immer\\
\textbf{Satz 1.5} & Sattelpunkt $\iff$ Gleichheit\\
\textbf{Dominanz} & strikt streichen ist unbedenklich\\
\bottomrule
\end{tabular}
\end{center}

\begin{center}
\begin{minipage}{0.88\textwidth}
\kursiv{Zwei der vier Matrizen aus Aufgabe 1.1 hatten keinen Sattelpunkt. Das nächste Kapitel zeigt, dass die Lücke stets verschwindet, sobald man eine einzige Zutat erlaubt: den Zufall.}
\end{minipage}
\end{center}
\end{merksatz}
